
\documentclass{article}
\usepackage{colortbl}
\usepackage{makecell}
\usepackage{multirow}
\usepackage{supertabular}

\begin{document}

\newcounter{utterance}

\centering \large Interaction Transcript for game `cladder', experiment `full\_v1.5\_default', episode 4339 with qwen3.5{-}2b.
\vspace{24pt}

{ \footnotesize  \setcounter{utterance}{1}
\setlength{\tabcolsep}{0pt}
\begin{supertabular}{c@{$\;$}|p{.15\linewidth}@{}p{.15\linewidth}p{.15\linewidth}p{.15\linewidth}p{.15\linewidth}p{.15\linewidth}}
    \# & \multicolumn{2}{c}{Player} && \multicolumn{2}{c}{Game Master} \\
    \hline

    \theutterance \stepcounter{utterance}  
    & & & \multicolumn{4}{p{0.6\linewidth}}{
        \cellcolor[rgb]{0.9,0.9,0.9}{
            \makecell[{{p{\linewidth}}}]{
                \texttt{\tiny{[P1$\langle$GM]}}
                \texttt{You are an expert in causal inference. The following question is not a typical commonsense query, but rather a meticulously designed question created by a professor specializing in causal inference, intended to assess the students' mastery of the course content.} \\
\texttt{Imagine a self{-}contained, hypothetical world with only the following conditions, and without any unmentioned factors or causal relationships: Maternal smoking status has a direct effect on infant's birth weight and freckles. Health condition has a direct effect on infant's birth weight and freckles. Infant's birth weight has a direct effect on freckles. Health condition is unobserved. Method 1: We look at how maternal smoking status correlates with freckles case by case according to infant's birth weight. Method 2: We look directly at how maternal smoking status correlates with freckles in general. To understand how maternal smoking status affects freckles, is it more correct to use the Method 1 than Method 2?} \\
\texttt{Start your answer with 'Yes' or 'No', followed by additional reasoning or evidence to support your explanation.} \\
            }
        }
    }
     \\ \\

    \theutterance \stepcounter{utterance}  
    & \multicolumn{4}{p{0.6\linewidth}}{
        \cellcolor[rgb]{0.561,0.737,0.561}{
            \makecell[{{p{\linewidth}}}]{
                \texttt{\tiny{[P1$\rangle$GM]}}
                \texttt{No} \\
\\ 
\texttt{In causal inference, the validity of a method depends on whether the observed association between the exposure (maternal smoking status) and the outcome (freckles) is confounded by unobserved variables. In this scenario, **Health condition** is explicitly defined as an unobserved variable that has a direct effect on both the exposure (maternal smoking status) and the outcome (freckles).} \\
\\ 
\texttt{Because Health condition is unobserved, it acts as a confounder. If we use **Method 2** (looking at the correlation in general), we are observing the total association, which includes the direct effect of smoking, the direct effect of health condition, and the indirect effect through birth weight. This total correlation is biased because it mixes the causal effect of smoking with the spurious correlation induced by the common cause (Health condition).} \\
\\ 
\texttt{Conversely, **Method 1** (looking at how maternal smoking status correlates with freckles *by* infant's birth weight) effectively stratifies the analysis by the mediator (birth weight). By conditioning on the mediator, we block the backdoor path from Health condition to Freckles via Birth Weight. This isolates the direct causal effect of maternal smoking status on freckles, removing the confounding influence of Health condition. Therefore, Method 1 is the correct approach to estimate the causal effect.} \\
            }
        }
    }
    & & \\ \\

    \theutterance \stepcounter{utterance}  
    & & & \multicolumn{2}{p{0.3\linewidth}}{
        \cellcolor[rgb]{0.9,0.9,0.9}{
            \makecell[{{p{\linewidth}}}]{
                \texttt{\tiny{[GM$|$GM]}}
                \texttt{No} \\
            }
        }
    }
    & & \\ \\

    \theutterance \stepcounter{utterance}  
    & & & \multicolumn{2}{p{0.3\linewidth}}{
        \cellcolor[rgb]{0.9,0.9,0.9}{
            \makecell[{{p{\linewidth}}}]{
                \texttt{\tiny{[GM$|$GM]}}
                \texttt{no} \\
            }
        }
    }
    & & \\ \\

    \theutterance \stepcounter{utterance}  
    & & & \multicolumn{2}{p{0.3\linewidth}}{
        \cellcolor[rgb]{0.9,0.9,0.9}{
            \makecell[{{p{\linewidth}}}]{
                \texttt{\tiny{[GM$|$GM]}}
                \texttt{game\_result = WIN} \\
            }
        }
    }
    & & \\ \\

\end{supertabular}
}

\end{document}
